%!TEX root = main.tex

\begin{itemize}
\item $\typeofexp(x, tmap)$ for name $x$ in $M$ is defined as follows. 
\begin{itemize}
  \item If $tmap(id(x)) = (\uintimp, n)$ or $(\sintimp, n)$ for some $n \in \Nat$, then the $\typeofexp(x, tmap)$ is defined as $(\uintimp, n')$ or $(\sintimp, n')$ where $n' \in \Nat$ is nondeterministically chosen. This nondeterministic choice is safe since expressions are used on the right-hand side of connect statements and all the chosen widths have to satisfy the {\typecompatible} constraints. 
%
\item If $tmap(id(x)) \in \gtypes \setminus \{(\uintimp, n), (\sintimp, n) \mid n \in \Nat\}$, then $\typeofexp(x, tmap) = tmap(id(x))$.
%
\item If $tmap(id(x))$ is an aggregate type, then {\typeofexp} inductively make nondeterministic choices for all the ground types of implicit widths in it. 
\end{itemize}
%
\item $\typeofexp(x.f, tmap)$ 
  is inductively defined as follows: Suppose $\typeofexp(x, tmap) = (\btype, ft)$. If $f$ occurs in $ft$, that is, $ft$ contains a triple $(f, b, t)$, then $\typeofexp(x.f, tmap)$ is defined as $t$, otherwise, $\typeofexp(x.f, tmap) = \svnone$. 
%
\item $\typeofexp(r[n], tmap)$ is inductively defined as follows: Suppose $\typeofexp(r, tmap) = (\vtype, t[m])$ for some $m \in \Nat$. If $n < m$, then $\typeofexp(r[n], tmap) = t$, otherwise, $\typeofexp(r[n], tmap) = \svnone$. 
%
\item $\typeofexp(\mbox{\prettyll{mux}}(c, e_1,e_2), tmap)= t$ is defined as follows: If $\typeofexp(c, tmap)= (\uint, 1)$ or $(\uintimp, 1)$, and $t_1 = \typeofexp(e_1, tmap)$ is equivalent to $t_2 = \typeofexp(e_2, tmap)$ after ignoring the widths (e.g. $\uint$ is equivalent to $\uint$, but not $\sint$), then $t$ is obtained from $t_1, t_2$ by taking the maximum of the widths. 
%
\item $\typeofexp(uop\ e_1)$ and $\typeofexp(e_1 \ bop\ e_2)$ for unary operations $uop$ and binary operations $bop$ can be defined according to the semantics of these operations. We  omit their definitions here, to avoid tediousness.
\end{itemize}
